\subsection{Physical-name density and finite-age coherence}\label{v:density-section}
The compiler of Section~\ref{sec:compilation}, including Section~\ref{s:probe}, gives
independent physical site-round names \(H=A\lor E\), with
\(\mu=\Pr(E)\le p=\Pr(H)\le2\mu\). A name may corrupt both its old
copies and its new output. The following supplies the precise density
estimate needed here, including that double appearance.

\begin{lemma}[All-path density]\label{v:density}
Let the names \(H(z,t)\) be independent with \(\Pr(H(z,t))\le p\).
Let \(\mathcal M_{\rm raw}\) be the least finite-start field that contains
every seed and marks a vertex whose marked ordinary parents
\((u+j,P+s,t-1)\), \(j\in\{-1,0,1\}^2\), \(s\) in the three replica boxes
\eqref{eq:c-boxes} acting on \((P_1,P_2)\), lie in two distinct boxes.
Seed \((z,t)\) when \(H(z,t)\) or \(H(z,t+1)\) occurs.
For an integer \(\ell\ge1\) let \(\mathcal D_\ell(z,t)\) mean that some
ordinary backward \(\ell\)-edge path from \((z,t)\) has a marked ordinary
parent at no fewer than \(\ell/1404\) of its first \(\ell\) vertices. Put
\[
 \Delta_{\rm raw}=13\,824,\quad B_{\rm raw}=2\Delta_{\rm raw}^2=382\,205\,952,
 \qquad p_{\rm dens}=B_{\rm raw}^{-67\,440}/64.
\]
For \(p\le p_{\rm dens}\),
\begin{equation}\label{v:density-bound}
 \Pr(\mathcal D_\ell(z,t))\le
 4\sum_{f\ge\lceil\ell/16\,848\rceil}
 [2B_{\rm raw}^{16\,860}\sqrt p]^f\le\frac{16}{3}e^{-\ell}.
\end{equation}
This holds simultaneously over all paths, including paths selected
after all innovations and initial words are known, for finite-start
marks without free initial seeds, and in injectively unwrapped torus
cones. Replacement values and initial words are unrestricted.
\end{lemma}
\par\noindent\emph{Proof idea.} Attach a complete three-charge contour to each selected marked parent and restart below its terminal time. The resulting graph controls every path in one count, with a separate payment for the two appearances of a physical source.

\begin{proof}
Use physical potentials
\[
 \psi_0=3(P_1+P_2)+350t,\quad
 \psi_1=-3P_1-175t,\quad \psi_2=-3P_2-175t.
\]
They sum to zero. Every distinct-box support pair offers each charge a
backward edge gaining at least 130; arbitrary loss is at most 350.
Every mark has a finite acyclic derivation ending at seeds.

Label frontier vertices by their descendant terminal sets, joining
vertices whose sets overlap. Expose the latest frontier vertex and its
two supports. An outside overlap component cannot be joined; the current
component splits into at most two daughters, each containing one of
these supports. Start three good arms at the root, one of each charge,
and retain one endpoint per charge per component. At a split add a
three-charge source at the exposed vertex, supplying missing endpoints;
only its first edges may be bad. Frontier vertices with no arm endpoint
must also be retained: they can carry the support through which a
split daughter is reached in one edge. Such exposures can create sources
in disconnected arm components; retain the component of the root source,
in which every source emits, and every terminal receives, exactly one
arm of each charge. If \(f'\) is its terminal count, each charge thus
bijects \(f'\) sources and \(f'\) terminals, with \(f'-1\) added sources.
Each bad first edge costs at most \(350+130=480\) relative to a good
edge, so balanced cancellation, with arm multiplicity, gives
\[
 130e_{\rm arm}\le480\cdot3(f'-1),\qquad e_{\rm arm}\le12(f'-1).
\]
A seeded root has one terminal and zero edges. Projection reduces size.
A minimum-time vertex is a terminal, and the root is the maximum, so the
whole connected contour has time span at most its edge budget.

Take the latest remaining marked-parent occurrence on a density path,
add its one-edge spur and its full contour, and skip occurrences whose
parent time is at or above the earliest terminal time of that contour.
Repeat, retaining the final contour whole even below the path endpoint.
The selected formal terminal slabs are disjoint. A contour with \(f_i\)
terminals covers at most \(12(f_i-1)+1\le12f_i\) positions, whence
\(f=\sum f_i\ge\lceil\ell/16\,848\rceil\). The entire path, all spurs
and all contours form a connected graph with at most
\(\ell+12f\le16\,860f\) edges. Its undirected degree is at most
\(2\cdot9\cdot768=13\,824\). Doubled-edge traversals and terminal subsets
have at most \(2B_{\rm raw}^m\) descriptions for \(m\) edges; summing
\(m\le16\,860f\) costs at most \(4B_{\rm raw}^{16\,860f}\). This counts
the whole path, not only the attached contours.

For \(f\) distinct formal terminals choose one of their two possible
physical names. There are at most \(2^f\) choices. A physical name seeds
at most two formal vertices, so each choice uses at least
\(\lceil f/2\rceil\) distinct independent names. Thus the terminal
probability is at most \(2^fp^{\lceil f/2\rceil}\le2^fp^{f/2}\).
Apply this once to the complete terminal union: adjacent formal slabs
can share an old/new name. This proves the series in
\eqref{v:density-bound}. At the stated rate bound its base is at most
\(B_{\rm raw}^{-16\,860}/4\); since
\(16\,860\log B_{\rm raw}/16\,848>1\), the series is at most
\(16e^{-\ell}/3\). Finite-start closure only deletes seeds; injective
cones preserve their law.
\end{proof}

The contour argument is that of \cite[Lemmas~12.3 and~12.4]{DMFixed2026}.
For each compiled source name, which has an old and a new appearance, this recovers equation (L.1) of \cite[Section~10.2]{DMNoise2026}. That equation was stated for single-vertex replacement marks. The bound also gives an explicit rate for the old/new density estimate of Lemma~\ref{c:coherence}. Its exact \(e^{-\ell}\) tail is stronger
than the downstream polynomial unions require. At the weaker cap
\(B_{\rm raw}^{-33\,720}/64\), the same count gives only
\((16/3)\exp[-(\log4)\ell/16\,848]\); the windows of Section~\ref{v:parameters} lie below the stronger cap.

The pairing argument of Lemma~\ref{c:coherence} (compare
\cite[Lemma~10.2]{DMNoise2026}) carries over to Moore inputs: two
unmarked nonbottom full-core readings at the same \((u,t)\) and
transverse distance \(d\) agree outside their two \(\mathcal D_\ell\)
events when \(t\ge\ell\ge\max(1,2d)\). Indeed at an unmarked root at
most one box has marked inputs, leaving at least \(512-256\) clean
successful replicas. Away from marked-parent positions a signed unit
shift gives \(3\cdot16\cdot15=720\) replica pairs. At most
\(2(768-512)=512\) fail, leaving 208 distance-reducing pairs.
Different outputs of the same deterministic Moore function force a
common differing input index. Follow that index at both roots. At an
exceptional step choose any clean successful replicas; distance rises
by at most 350. After the first step the readings are actual first
coordinates, so even time-zero distance-zero conflict is impossible.
With fewer than \(2\ell/1404\) exceptional paired steps the final
distance is less than \(d-\ell/2\), a contradiction. All-parent bottom
spreading ensures all followed inputs are nonbottom. This proves the
finite-age comparison for the core and the ordinary second proposal.

We will also use the physical source slabs of the proofs of
Lemma~\ref{c:coherence} and Propositions~\ref{c:late-live}
and~\ref{c:path-tail}. A complete witness reserves
all its old/new names; restart strictly below its earliest physical
name minus two. At a boundary formal time \(b\), reserve names through
\(b+1\) and their appearances in \(\{b,b+1\}\); future names begin at
\(b+2\). The last witness is kept whole. These conventions separate
physical supports, including names whose two formal appearances straddle
a cut. They are used both for deep historical queries and for rare
contours along quantum paths.
